\documentclass[10pt,a4paper]{amsart}
\usepackage[margin=19mm]{geometry}
\usepackage{amsmath,amssymb,mathtools}
\usepackage[scr=rsfs,cal=euler]{mathalfa}
\usepackage{xfrac}
\usepackage[unicode,hidelinks,bookmarks=false]{hyperref}
\newcommand{\ie}{i.e.\ }
\newcommand{\cf}{cf.\ }
\newcommand{\resp}{respectively\ }
\newcommand{\Subsection}{\subsection}
\providecommand{\nobreakdash}{\nobreak\hbox{-}}
\setlength{\emergencystretch}{2em}
\multlinegap0pt
\usepackage{bm}
\usepackage{xcolor}
\usepackage{amssymb}
\newtheorem{theorem}{Theorem}
\newcommand{\bydef}{\stackrel{\rm def}{=}}
\title{Holomorphic Interpolation of Multivariate Completely Monotone Functions}
\author{Mainak Bhowmik, Agniva Chatterjee, Mihai Putinar}
\date{Source: arXiv:2606.12102v1 (CC BY 4.0)}
\begin{document}
\maketitle

\noindent\emph{Source and licence.} Holomorphic Interpolation of Multivariate Completely Monotone Functions. Version arXiv:2606.12102v1 (CC BY 4.0), distributed by arXiv under the Creative Commons Attribution 4.0 International licence, http://creativecommons.org/licenses/by/4.0/, as recorded in the arXiv licence metadata for this exact version and shown on the abstract page https://arxiv.org/abs/2606.12102v1. Full licence notice: \texttt{licenses/arxiv-2606.12102-CC-BY-4.0.txt}.\par\smallskip

\noindent\textit{Editorial scope.} Counted unit: the article's theorem conv (source0.tex lines 1114--1125). The article's complete original proof of it is delivered with it (source0.tex lines 1127--1332, one \texttt{proof} environment of 206 lines). No mathematics is written, repaired or re-derived: the statement and the proof are the article's own lines, in the article's own notation, and no retained line is substituted or re-flowed.  No further source-local context is needed: every cross-reference of every retained line resolves inside the two delivered spans.  Nothing was omitted from any retained span, so the package carries no omission record.  No retained line contains a citation, so the wrapper carries no bibliography and no bibliography entry.  No retained line contains a figure environment, an image-inclusion command or drawing code, so no image is delivered and the package declares no asset dependency.  Editorial formatting only, in addition: an AMS \texttt{amsart} preamble that restates the article's own declarations and macros used by the retained lines, namely the environments \texttt{theorem} on separate counters, together with the standard packages those lines need.  The retained environments print in this document as Theorem 1 in order of appearance, whereas the article prints the same items with the numbering of its own full document.  The complete native source of the article as distributed in the arXiv e-print for this exact version is bundled unchanged as \texttt{sources/202599/source0.tex}.
	\begin{theorem}\label{th:Lap conv}
		Let $\Gamma \subset {\mathbb R}^d$ be an acute closed convex solid cone and let $\mu$ be a positive measure of moderate growth defined on $\Gamma^\ast$.
		Fix a point $\lambda \in \operatorname{int} \Gamma$. For every $n \in \mathbb{N}$, there exists a Wigner distribution associated 
		to a $d$-tuple of real symmetric matrices, such that its Laplace transform $F_n(z)$ is matching the Taylor coefficients of
		the Laplace transform
		$$ F(z) = \int_{\Gamma^\ast} \exp{(-z\cdot x)}\ d\mu(x), \ \ \Re z \in \operatorname{int} \Gamma,$$ 
        up to order $2n+1$ at $z = \lambda$. 
        
		The functions $F_n$ are entire of exponential type, and they are directional exponential sums on the affine cone $\lambda+\Gamma$. 
        
        Moreover, the sequence $(F_n)_n$ converges to $F$ on the open tube domain $T_{\lambda + \Gamma}$ associated with the cone $\lambda + \Gamma$, in the Fr\'echet topology of uniform convergence on compact subsets.
	\end{theorem}

\begin{proof}
Let us start with the Laplace transform $F$ of a positive measure $\mu$ having moderate growth and supported in $\Gamma^*$ and a point $\lambda \in \operatorname{int} \Gamma$ as in the statement of the theorem. Fix $n\in \mathbb{N}$.

	Let $H_n$ denote the finite dimensional subspace of $L_{\mathbb{R}}^2(\Gamma^*, \mu)$ defined as
    \begin{align*}
        H_n= \operatorname{span}_\mathbb{R} \{x^\alpha e^{-\frac{\lambda}{2} \cdot x} : \ \alpha \in \mathbb{N}_0^d, \ |\alpha| \leq n\}.
        \end{align*}
    
    As in the preceding section, for every $j \in \{1, \dots, d\}$, we consider the compression: 
    $$
    X^{(n)}_j = P_n X_j P_n
    $$ 
    of $X_j$, the multiplication by the variable $x_j$ on $L^2(\Gamma^*, \mu)$, on $H_n$. Here $P_n$ denotes the orthogonal projection onto $H_n$. Note that $X_j$ might be an unbounded, closed graph operator.
	
	The tuple $\boldsymbol{X}^{(n)} = (X^{(n)}_1, \dots, X^{(n)}_d)$ of real symmetric, finite dimensional space operators has joint numerical range $W(\boldsymbol{X}^{(n)})$ contained in the closed convex cone $\Gamma^\ast$. Indeed, the convexity of the closed  cone $\Gamma^\ast$ yields
	$$ \langle \boldsymbol{X}^{(n)} \psi, \psi \rangle =\left(\int_{\Gamma^\ast} x_1 \psi(x)^2 d\mu(x), \ldots, \int_{\Gamma^\ast} x_d \psi(x)^2 d\mu(x) \right) \in \Gamma^\ast,$$
	whenever $\psi(x) = \sum_{|\alpha| \leq n} c_\alpha x^\alpha e^{-\frac{\lambda}{2} \cdot x}$ in $H_n$ satisfies 
	$ \|\psi\| = 1.$
	Accordingly,
	$$ \Re (p\cdot \boldsymbol{X}^{(n)}) \geq 0$$ as an operator, for every $ p \in \Gamma$. 
We define the associated entire function,
	$$ F_n(z) \bydef \langle \exp (- (z- \lambda) \cdot \boldsymbol{X}^{(n)}) e^{-\frac{\lambda}{2} \cdot x}, e^{-\frac{\lambda}{2} \cdot x} \rangle, \quad z\in \mathbb{C}^d.$$
	Denote $u(x) =  e^{-\frac{\lambda}{2} .\cdot x}$ and let ${\mathcal W}_u(\boldsymbol{X}^{(n)})$ be the Wigner distribution associated to $\boldsymbol{X}^{(n)}$ and vector $u$ such that
	$$( {\mathcal W}_u(\boldsymbol{X}^{(n)})(\xi), \phi(x \cdot \xi)) = \langle \phi( x\cdot \boldsymbol{X}^{(n)}) u, u\rangle, \ \ \phi \in {\mathcal E}({\mathbb R}).$$
	Consequently, 
	$$ F_n(z + \lambda) = ({\mathcal W}_u(\boldsymbol{X}^{(n)})(\xi),\,  e^{-z \cdot \xi})$$
	is the Laplace transform of the distribution ${\mathcal W}_u(\boldsymbol{X}^{(n)})$. Clearly, $F_n(z)$ is entire function of exponential type. 

\vspace{2mm}

  \noindent \textbf{Matching  the Taylor coefficients:} We now show that the partial derivatives in the variables $z$ of $F_n$ and $F$ are identical at $z = \lambda$, up to order $2n+1$.
    
For $\alpha=(\alpha_1, \dots, \alpha_d) \in \mathbb{N}_0^d$, the coefficient of $\frac{1}{|\alpha|!}(z-\lambda)^\alpha$ in the non-commutative power series expansion of $\exp (- (z- \lambda) \cdot \boldsymbol{X}^{(n)})$ at $\lambda$, turns out to be the sum of the words in $-X^{(n)}_1, \dots,-X^{(n)}_d$ containing $\alpha_j$ many $X^{(n)}_j$ for each $j$. We denote the collection of such words by $\mathcal{M}(\alpha)$.
Thus, for $|\alpha| \leq 2n+1$, the coefficient of $(z-\lambda)^\alpha$ in the power series of $F_n(z)$ at $\lambda$
would become:
\begin{align}\label{Eq:NC-sum}
\frac{1}{|\alpha|!} \sum_{M \in \mathcal{M}(\alpha)} \langle M e^{-\frac{\lambda}{2} \cdot x}, e^{-\frac{\lambda}{2} \cdot x} \rangle.
\end{align}

Note that, if $|\alpha| \leq n$, then for any word $M \in \mathcal{M}(\alpha)$,
\begin{align}\label{Eq:NC-Laplace}
M e^{-\frac{\lambda}{2} \cdot x} = (-1)^{|\alpha|} [\boldsymbol{X}^{(n)}]^\alpha e^{-\frac{\lambda}{2} \cdot x}.
\end{align}
Moreover, if $|\alpha|\leq 2n+1$ and $M\in \mathcal{M}(\alpha)$, then we write $M=M_1M_2$ such that $M_1$ has length at most $n+1$ and $M_2$ has length at most $n$. So, each summand in \eqref{Eq:NC-sum} can be expressed in the form $\langle M_1 e^{-\frac{\lambda}{2} \cdot x}, M_2 e^{-\frac{\lambda}{2} \cdot x}\rangle$. Suppose, $M_j$ contains $s_{j, k}$ many $X^{(n)}_k$ in its expression as a word with $1\leq k \leq d$ and $\beta^{(j)} = (s_{j, 1}, \cdots, s_{j, d} )$ for $j=1,2$. Thus, using \eqref{Eq:NC-Laplace} for $M_1$ and $M_2$ we obtain 
\begin{align*}
\langle M_1 e^{-\frac{\lambda}{2} \cdot x}, M_2 e^{-\frac{\lambda}{2} \cdot x}\rangle &= \langle (-1)^{|\beta^{(1)}|} [\boldsymbol{X}^{(n)}]^{\beta^{(1)}} e^{-\frac{\lambda}{2} \cdot x}, (-1)^{|\beta^{(2)}|} [\boldsymbol{X}^{(n)}]^{\beta^{(2)}} e^{-\frac{\lambda}{2} \cdot x}\rangle \\
& =(-1)^{|\alpha|} \int_{\Gamma^*} x^\alpha e^{-\lambda \cdot x} d\mu(x).
\end{align*}
Therefore, \eqref{Eq:NC-sum} yields that the coefficient of $(z-\lambda)^\alpha$ in the power series of $F_n(z)$ at $z=\lambda$, for $|\alpha| \leq 2n+1$, is
$$
(-1)^{|\alpha|} \frac{|\mathcal{M}(\alpha)|}{|\alpha|!}\int_{\Gamma^*} x^\alpha e^{-\lambda \cdot x} d\mu(x) =  \frac{(-1)^{|\alpha|}}{\alpha !} \int_{\Gamma^*} x^\alpha e^{-\lambda \cdot x} d\mu(x).
$$
Hence, for $|\alpha| \leq 2n+1$
$$
\partial^\alpha_ z F_n |_{z=\lambda} = \partial^\alpha_ z F |_{z=\lambda}, 
$$
i.e., $F_n$ and $F$ have the matching Taylor coefficients up to order $2n+1$ at $\lambda$. 

\vspace{2mm}

 \noindent \textbf{Directional exponential sum:} Denote the function $e^{-\frac{\lambda}{2} \cdot x}$ by $u$. For a fixed $\omega \in \Gamma$, by orthogonal diagonalization of the real symmetric operator $ \omega \cdot X^{(n)}$ we obtain that
	\begin{align}\label{Eq:dir-exp-sum}
	   F_n(\lambda + t \omega) = \sum_{k=1}^N |c_k|^2 \exp{(-t \gamma_k )}, \ \ t\in [0, +\infty) 
	\end{align}
    is an sum of $N ( \leq \dim H_n)$ exponentials where the nodes $\gamma_k \in \sigma (\omega \cdot \boldsymbol{X}^{(n)}) \subset [0, +\infty)$ and the coefficients $|c_k|^2 \geq 0$ depend on the direction $\omega$. Therefore, $F_n$ is a directional exponential sum on the affine cone $\lambda +\Gamma$.
    In fact, if $U$ is an orthogonal matrix with an ordered orthonormal eigen-basis on $H_n$ as its column, then the real numbers $c_1, \dots, c_d$ are the inner products of the ordered basis vectors with $u$, respectively. So, $\sum_{j=1}^N |c_j|^2 = \|u\|^2 = F(\lambda)$. 
    
     Note that $F_{n, \omega}$ has a unique holomorphic extension to $\mathbb{H}^+_R = [0,\infty) + i \mathbb{R}$ given by
    $$
    F_{n, \omega}(z) =  \langle \exp(- z \omega \cdot \boldsymbol{X}^{(n)})u, u \rangle, \ \ z\in \mathbb{H}^+_R.
    $$
    Further, $\omega \cdot \boldsymbol{X}^{(n)}$ being positive semi-definite,
    $$ \| \exp(-z  \omega \cdot \boldsymbol{X}^{(n)}) \| \leq 1 \quad \text{for } z\in \mathbb{H}^+_R.$$
    
    Also, from \eqref{Eq:dir-exp-sum}, we see that, for $t+is \in \mathbb{H}^+_R$,
    $$
    |F_{n, \omega}(t +is)| =\left| \sum_{k=1}^N |c_k|^2 \exp{(-(t+is) \gamma_k )} \right| \leq \sum_{k=1}^N |c_k|^2 = F(\lambda).
    $$
    Further more, for $1\leq k \leq n$, we have
    \begin{align*}
        \frac{d^k F_{n, \omega}}{dz^k}(z) = (-1)^k \langle (\omega \cdot \boldsymbol{X}^{(n)})^k \exp(-z \ \omega \cdot \boldsymbol{X}^{(n)})) u, u \rangle.
    \end{align*}
For $1\leq k \le n$, the vector $(\omega \cdot \boldsymbol{X}^{(n)})^k u$ in $H_n$ can be expressed as 
$$
(\omega \cdot \boldsymbol{X}^{(n)})^k u = \sum_{\alpha \in \mathbb{N}_0^d; |\alpha|=k} \frac{k!}{\alpha!}\omega^\alpha [\boldsymbol{X}^{(n)}]^\alpha u.
$$
This is because, for any word $M$ of length $k$ $(\leq n)$ with $\alpha_j$ many $X^{(n)}_j$ as its letters for $1\leq j \leq d$, $M u = [\boldsymbol{X}^{(n)}]^\alpha u$ as illustrated in \eqref{Eq:NC-Laplace}. Therefore,
\begin{align*}
 \| (\omega \cdot \boldsymbol{X}^{(n)})^k u  \|^2 &= \sum_{\alpha, \beta \in \mathbb{N}_0^d; |\alpha|=|\beta|=k} \frac{(k!)^2}{\alpha! \beta!} \omega^{\alpha + \beta}\langle [\boldsymbol{X}^{(n)}]^\alpha u,  [\boldsymbol{X}^{(n)}]^\beta u\rangle \\
 &= \sum_{\alpha, \beta \in \mathbb{N}_0^d; |\alpha|=|\beta|=k} \frac{(k!)^2}{\alpha! \beta!} \omega^{\alpha + \beta} \int_{\Gamma^*} x^{\alpha + \beta} e^{-\lambda \cdot x} d\mu(x)\\
 &= \sum_{\alpha, \beta \in \mathbb{N}_0^d; |\alpha|=|\beta|=k} \frac{(k!)^2}{\alpha! \beta!} \omega^{\alpha + \beta} \partial^{\alpha + \beta}_z F|_{z=\lambda}\\
 &=\frac{d^k F_\omega}{dt^k}(0),
\end{align*}
    where $F_\omega$ is the slice function given by 
    $$
    F_\omega(t) \bydef F(\lambda + t \omega), \quad t\in [0, \infty).
    $$
    Whenever $0\leq k \leq n$, we assemble these information to get
    \begin{align*}
        \left| \frac{d^k F_{n, \omega}}{dz^k}(z) \right| &= \left| \langle  \exp(-z \ \omega \cdot \boldsymbol{X}^{(n)})) u, (\omega \cdot \boldsymbol{X}^{(n)})^k u \rangle \right| \\
        & \leq \|\exp(-z \ \omega \cdot \boldsymbol{X}^{(n)})) u\| \| (\omega \cdot \boldsymbol{X}^{(n)})^k u\| \\
        & \leq (F(\lambda))^{1/2} \left(\frac{d^k F_\omega}{d t^k}(0)\right)^{1/2}.
    \end{align*}


Thus, for each $k \in \mathbb{N}_0$, the sequence of smooth functions $\{F_{n, \omega}\}_n$ on $\overline{\mathbb{H}^+_U}$ satisfies
\begin{align} \label{Eq:Arzela-cond}
    \left| \frac{d^k F_{n, \omega}}{dz^k}(z) \right|  \leq (F(\lambda))^{1/2} \left(\frac{d^k F_\omega}{d t^k}(0)\right)^{1/2}, \ \text{for all } n \geq k. 
\end{align}

\vspace{2mm}
    
\noindent \textbf{Convergence of the sequence of approximants:} We complete the proof of the theorem by showing that the sequence of entire function $\{F_n\}_n$ converges to $F$ in $\mathcal{O}(T_{\lambda + \Gamma})$, where $T_{\lambda + \Gamma} = (\lambda + \operatorname{int} \Gamma) + i \ \mathbb{R}^d$.

\vspace{2mm}

\noindent \textbf{Claim 1:} The sequence $\{F_n\}_n$ is uniformly bounded on $T_{\lambda + \Gamma}$.

\noindent \textbf{Proof of Claim 1:} Take any $z= \lambda + p + i q \in T_{\lambda + \Gamma}$ where $p \in \operatorname{int} \Gamma$ and $q\in \mathbb{R}^d$. Then,
$$
F_n(z) = \langle \exp(-P - i Q) u, u \rangle,
$$
where $P= p \cdot \boldsymbol{X}^{(n)}$ and $Q= q\cdot \boldsymbol{X}^{(n)}$ which are both real symmetric operators on $H_n$. Define the $H_n$-valued function
   $$
   v(t) \bydef \exp(-t(P+ iQ)) u, \quad t\in [0, \infty).
   $$
Then 
$$
v'(t) \bydef \frac{d v}{dt}(t) = -(P+i Q) v(t).
$$ and 
\begin{align} \label{Eq:dervative}
    \frac{d}{dt}(\|v(t)\|^2) = \langle v'(t),  v(t) \rangle + \langle v(t), v'(t) \rangle = -2\langle  P \ v(t),  v(t)\rangle.
\end{align}
Denoting $\hat{v}(t) = \frac{v(t)}{\|v(t)\|}$, we have 
\begin{align*}
    \langle  P \ v(t),  v(t)\rangle = \|v(t)\|^2 \langle  p \cdot \boldsymbol{X}^{(n)} \ \hat{v}(t),  \hat{v}(t)\rangle = \|v(t)\|^2  (p \cdot \langle \boldsymbol{X}^{(n)} \ \hat{v}(t),  \hat{v}(t)\rangle) \geq 0, 
\end{align*}
since $\langle \boldsymbol{X}^{(n)} \ \hat{v}(t),  \hat{v}(t)\rangle \in W(\boldsymbol{X}^{(n)}) \subseteq \Gamma^*$. Thus, from \eqref{Eq:dervative} we have
$$
 \frac{d}{dt}(\|v(t)\|^2) \leq 0 \quad \text{for } t\in [0, \infty), 
$$
and so, for $t\in [0, \infty)$,
$$
\|v(t)\|^2 \leq \|v(0)\|^2 = \|u\|^2 = F(\lambda).
$$
Putting $t=1$, we have 
$$
|F_n(z)| \leq \|\exp(-(P+iQ))u \| \|u\| \leq \|u\|^2 = F(\lambda),
$$
for any $z \in T_{\lambda + \Gamma}$. This proves our Claim 1.

\vspace{2mm}

Now, by Montel's theorem we get a subsequence $\{F_{n_j}\}_j$ of $\{F_n\}_n$ such that $F_{n_j}$ converges to $G$ in the Fr\'echet topology of $\mathcal{O}(T_{\lambda+\Gamma})$. We shall prove that $G=F$ in $T_{\lambda + \Gamma}$. Then the original sequence $\{F_n\}_n$ itself will converge to $F$ in $\mathcal{O}(T_{\lambda+\Gamma})$. Therefore, without any loss of generality, we assume that $F_n$ converges to $G$. 

At this stage, it might be tempting to conclude that $G=F$ follows from the standard normal family argument as we have matching derivatives of $F_n$ and $F$ up to order $2n+1$ at the point $\lambda$,. But this is not the case because $\lambda$ is not an interior point of $T_{\lambda +\Gamma}$. So, we analyze the slice functions along each direction in $\Gamma$.

Fix a direction $\omega \in \Gamma$. We know that the slice function $F_{n, \omega}$ satisfies the bounds in \eqref{Eq:Arzela-cond}. Applying Arzela-Ascoli theorem, we get a subsequence $\{F_{n_j, \omega}\}$ such that for each $k\in \mathbb{N}_0$,
\begin{align}\label{Eq:Arzela-Ascoli}
\frac{d^k F_{n_j, \omega}}{dz^k} \  \rightarrow \ \frac{d^k g_\omega} {dz^k}, \quad \text{as } j\to \infty,
\end{align}
uniformly on compact subsets of $\overline{\mathbb{H}^+_R}$, for some smooth function $g_\omega$ on $\overline{\mathbb{H}^+_R}$. In particular, $g_\omega \in \mathcal{O}(\mathbb{H}^+_R)$.

\vspace{2mm}

\noindent \textbf{Claim 2:} For each $k\in \mathbb{N}_0$,
\begin{align}\label{Eq:equal-derivatives}
    \frac{d^k F_\omega}{dt^k}(0) = \frac{d^k g_\omega}{dz^k}(0).
\end{align}

\noindent \textbf{Proof of Claim 2:} From \eqref{Eq:Arzela-Ascoli}, we have
\begin{align}\label{Eq:conv-zero}
    \frac{d^k F_{n_j, \omega}}{dz^k}(0) \  \rightarrow \ \frac{d^k g_\omega} {dz^k}(0), \quad \text{as } j\to \infty.
\end{align}
Again for $k \leq n_j$,
\begin{align*}
    \frac{d^k F_{n_j, \omega}}{dz^k}(0) &= (-1)^k \langle (\omega \cdot \boldsymbol{X}^{(n_j)} )^k u, u \rangle  \\
    &= (-1)^k \sum_{\alpha \in \mathbb{N}_0^d; |\alpha|=k} \frac{k!}{\alpha!}\omega^\alpha \langle  [\boldsymbol{X}^{(n_j)} ]^\alpha u, u \rangle \\
    &=(-1)^k \sum_{\alpha \in \mathbb{N}_0^d; |\alpha|=k} \frac{k!}{\alpha!}\omega^\alpha \int_{\Gamma^*} x^\alpha e^{-\lambda \cdot x} d\mu(x) \\
    &= \sum_{\alpha \in \mathbb{N}_0^d; |\alpha|=k} \frac{k!}{\alpha!}\omega^\alpha \partial^\alpha_z F_\omega|_{z=\lambda}
    = \frac{d^k F_\omega}{dt^k}(0).
\end{align*}
The second equality above is obtained by using \eqref{Eq:NC-Laplace}.
Thus, \eqref{Eq:conv-zero} yields the equality as in \eqref{Eq:equal-derivatives} for every $k\in \mathbb{N}_0$. This proves our Claim 2.


\vspace{2mm}
\noindent \textbf{Concluding part:} Since $\lambda$ is an interior point of $\Gamma$, there exists $\epsilon > 0$ such that $\lambda + t\omega \in \operatorname{int} \Gamma$ for $t\in (-\epsilon, \epsilon)$. Also, the function $F$ is defined in the tube domain $T_\Gamma$. Thus, the function
$$
F_\omega(t) = F(\lambda + t\omega) = \int_{\Gamma^*} \exp(-(\lambda+ t\omega)\cdot x) d\mu(x)
$$
is also defined for $t\in (-\epsilon, \infty)$. However, $F$ being holomorphic in $T_\Gamma$, $F_\omega$ is real analytic in $(-\epsilon, \infty)$. So, $F_\omega$ has a power series representation in neighbourhood $(-\epsilon', \epsilon')$ of $0$. But the equality of derivatives of $F_\omega$ and $g_\omega$ at $0$ as in \eqref{Eq:equal-derivatives} implies that the same power series provides a real analytic extension of $g_\omega$ in $(-\epsilon', \epsilon')$. Remember that $g_\omega$ is holomorphic in $\mathbb{H}^+_R$. Thus, $g_\omega$ is real analytic in $(-\epsilon', \infty)$ and so, by the identity theorem for real analytic functions, $F_\omega = g_\omega$ on $(-\epsilon', \infty)$.

Therefore, for any arbitrary point $\lambda + p \in \lambda + \operatorname{int} \Gamma$. Then 
$$
\lim_{j\to \infty} F_{n_j}(\lambda + p) = \lim_{j\to \infty} F_{n_j, p} (1) = g_p(1) = F_p(1)=F(\lambda +p).
$$
Again, recall that $F_n$ converges to $G$ in $\mathcal{O}(T_{\lambda + \Gamma})$. So, 
$$
\lim_{j\to \infty} F_{n_j}(\lambda + p) = G(\lambda + p).
$$
Thus, $G(\lambda + p)=F(\lambda + p)$ and hence, $G=F$ in $\lambda + \operatorname{int} \Gamma$.  

Finally, we note that $\Gamma$ being a solid cone, $\lambda + \operatorname{int} \Gamma$ is a non-empty open set in $\mathbb{R}^d$ and so, by the Cauchy-Riemann equations for holomorphic functions in the open tube $T_{\lambda + \Gamma}$ we see that $\lambda + \operatorname{int} \Gamma$ is a uniqueness set for $T_{\lambda +\Gamma}$. Hence, $F=G$ in $T_{\lambda + \Gamma}$. This completes the proof.

\end{proof}

\end{document}
